\subsection{معرفی و انگیزه}
این ماژول در چارچوب پژوهش حاضر پیاده‌سازی و ارزیابی شده و فاقد مقاله مستقل جداگانه در ادبیات تحقیق است.

مکانیزم توجه \lr{Hybrid Attention} (توجه ترکیبی چندشاخه با گیت اختصاصی به ازای هر سر) با هدف تلفیق قابلیت‌های بیان‌گری و تطبیق‌پذیری بالای مکانیزم توجه سنتی بر پایه تابع بیشینه هموار (\lr{Softmax Attention}) و کارایی محاسباتی خطی مکانیزم‌های بازگشتی یا هسته‌محور (\lr{Linear Attention}) در پردازش توالی‌های طویل طراحی شده است. توجه استاندارد دارای پیچیدگی زمانی و مکانی $\mathcal{O}(N^2)$ است که در طول دنباله‌های زیاد منجر به اشباع سریع حافظه کارت گرافیک (\lr{VRAM}) و افت محسوس نرخ پردازش توکن‌ها می‌گردد؛ از سوی دیگر، مکانیزم‌های توجه خطی با تقریب نگاشت هسته، هرچند حافظه و توان محاسباتی را به $\mathcal{O}(N \cdot d)$ تقلیل می‌دهند، اما در ثبت وابستگی‌های نحوی و ساختاری ریزدانه با کاهش ظرفیت مواجه می‌شوند.

نوآوری بنیادین این ماژول در پیاده‌سازی گیت تلفیق یادگیرنده برای هر سر توجه (\lr{Per-Head Learned Gate}) نهفته است. بر خلاف معماری‌هایی نظیر \lr{MEGA} که یک مقدار گیت سراسری مشترک را اعمال می‌کنند، در این ماژول به هر سر اجازه داده می‌شود به طور پویا و متناسب با بافت محتوایی ورودی، اولویت خود را میان شاخه دقیق دانس و شاخه خطی تعیین نماید. علاوه بر این، با راه‌اندازی اولیه بایاس گیت روی مقادیر بزرگ (\lr{bias = 10.0})، شبکه در آغاز فرآیند یادگیری کاملاً بر شاخه \lr{Softmax} منطبق بوده و با حداقل انحراف مقادیر اولیه‌ی مدل پایه، مسیر بهینه‌سازی را بدون افت کیفی اولیه طی می‌کند.
